% !TEX TS-program = pdflatex
% !TeX program = pdflatex
% !TEX encoding = UTF-8
% !TEX spellcheck = fr

\documentclass[xcolor=table]{beamer}


%\usepackage{fullpage}
%\usepackage[left=2.8cm,right=2.2cm,top=2 cm,bottom=2 cm]{geometry}
\setbeamersize{text margin left=10pt,text margin right=10pt}


\usepackage{amsmath,amssymb} 
\usepackage{fontenc}
\usepackage{textcomp}
\usepackage[utf8]{inputenc}
\usepackage{CJKutf8}
\usepackage[french]{babel}
\usepackage{arabtex,acolor}
\usepackage{txfonts}
\usepackage{tipa}
\usepackage[]{graphicx}
\usepackage{multirow}
\usepackage{hyperref}
\usepackage{colortbl}
\usepackage{tabularray}
\usepackage{listingsutf8}
\usepackage{wrapfig}
\usepackage{multicol}
\usepackage[export]{adjustbox} %for images in table, also for frame
\usepackage[many]{tcolorbox}
\usepackage{wasysym}
\usepackage[french,lined]{algorithm2e}
\usepackage{alltt} %verbatim with commands
\usepackage{longtable}
\usepackage{tabu}


\definecolor{lightblue}{HTML}{D0D2FF}
\definecolor{lightyellow}{HTML}{FFFFAA}
\definecolor{darkblue}{HTML}{0000BB}
\definecolor{olivegreen}{HTML}{006600}
\definecolor{darkgreen}{HTML}{008B45} %009B55
\definecolor{violet}{HTML}{6600CC}
%\definecolor{deeppink}{HTML}{FF1493}
\definecolor{orangey}{HTML}{FFBB00}

\usetheme{Warsaw} % Antibes Boadilla Warsaw

\beamertemplatenavigationsymbolsempty


\let\oldcite\cite
\renewcommand{\cite}[1]{{\bfseries\color{orangey}\oldcite{#1}}}

\tcbuselibrary{listings}

\newcommand{\kurl}[1]{{\scriptsize\bfseries\color{orangey}\url{#1}}}


%\renewcommand{\baselinestretch}{1.5}

\def\supit#1{\raisebox{0.8ex}{\small\it #1}\hspace{0.05em}}

\AtBeginSection{%
	\begin{frame}
		\sectionpage
	\end{frame}
}

\newcommand{\rottext}[2]{%
	\rotatebox{90}{%
	\begin{minipage}{#1}%
		\raggedleft#2%
	\end{minipage}%
	}%
}


\institute{ %
Laboratoire de la Communication dans les Systèmes Informatiques (LCSI)

\vspace{6pt}
École  nationale Supérieure d'Informatique (ESI, ex. INI), Alger, Algérie
}
\author[ \textbf{\footnotesize\insertframenumber/\inserttotalframenumber} \hspace*{1cm} ESI - ARIES Abdelkrime (Master 2022/2023)] %
{ARIES Abdelkrime}
%\titlegraphic{\includegraphics[height=1cm]{../img/esi-logo.png}%\hspace*{4.75cm}~
	
	
	\date{Master 2022/2023} %\today

\titlegraphic{%
	\includegraphics[height=1cm]{../img/esi-logo.png}%
	\hspace{2cm}
	\includegraphics[height=1cm]{../img/lcsi-logo.png}%
}


%\setbeamertemplate{headline}{}



\newcommand{\keyword}[1]{\textcolor{red}{\bfseries\itshape #1}}
\newcommand{\expword}[1]{\textcolor{olivegreen}{#1}}
\newcommand{\optword}[1]{\textcolor{violet}{\bfseries #1}}

\makeatletter
\newcommand\mysphere{%
	\parbox[t]{10pt}{\raisebox{0.2pt}{\beamer@usesphere{item projected}{bigsphere}}}}
\makeatother

%\let\oldtabular\tabular
%\let\endoldtabular\endtabular
%\renewenvironment{tabular}{\rowcolors{2}{white}{lightblue}\oldtabular\rowcolor{blue}}{\endoldtabular}


\NoAutoSpacing %french autospacing after ":"

\def\graphpath{}

\newcommand{\changegraphpath}[1]{\def\graphpath{#1}}


\newcommand{\vgraphpage}[2][.84\textheight]{%
%	\begin{center}%
		\includegraphics[height=#1]{\graphpath #2}%
%	\end{center}%
}

\newcommand{\hgraphpage}[2][\textwidth]{%
%	\begin{center}%
		\includegraphics[width=#1]{\graphpath #2}%
%	\end{center}%
}

\newcommand{\graphpage}[2][]{%
	\includegraphics[#1]{\graphpath #2}%
}

\bibliographystyle{apalike}

\newcommand{\insertbibliography}[2]{
	\appendix
	\section*{Bibliographie}
	\nocite{#2}
%	\makeatletter % to change template
%	\setbeamertemplate{headline}[default] % not mandatory, but I though it was better to set it blank
%	\def\beamer@entrycode{\vspace*{-\headheight}} % here is the part we are interested in :)
%	\makeatother
	\begin{multicols*}{2}[\frametitle{\insertsection} \usebeamertemplate*{frametitle}]%\usebeamertemplate*{frametitle}\frametitle{Références}
		\tiny
		\bibliography{#1}
	\end{multicols*}
}

\definecolor{my-grey}{RGB}{233, 233, 233}

\newcommand{\insertlicence}{
	\begin{frame}[plain]
	\frametitle{Licence : CC-BY 4.0}
%	\framesubtitle{Licence: CC-BY-NC 4.0}

	\begin{tcolorbox}[colback=cyan,
		colframe=cyan,  
		arc=0pt,outer arc=0pt,
		valign=top, 
		halign=center,
		width=\textwidth]
		
		\includegraphics[width=.5cm]{../img/licence/cc_icon_white_x2.png}
		\includegraphics[width=.5cm]{../img/licence/attribution_icon_white_x2.png}
		
		\color{white}
		\bfseries Attribution 4.0 International (CC BY 4.0) \\
		\tiny \url{https://creativecommons.org/licenses/by/4.0/deed.fr}
		
	\end{tcolorbox}\vspace{-.5cm}
	\begin{tcolorbox}[colback=my-grey,
		colframe=my-grey,  
		center, arc=0pt,outer arc=0pt,
		valign=top, 
		halign=left,
		width=\textwidth]
		
		\tiny
		
		\begin{center}
			\bfseries\Large
			Vous êtes autorisé à :
		\end{center}
		
		\begin{minipage}{0.83\textwidth}
			\begin{itemize}
				\item[] \textbf{Partager} — copier, distribuer et communiquer le matériel par tous moyens et sous tous formats
				\item[] \textbf{Adapter} — remixer, transformer et créer à partir du matériel
				pour toute utilisation, y compris commerciale.
			\end{itemize}
		\end{minipage}
		\begin{minipage}{0.15\textwidth}
			\includegraphics[width=\textwidth]{../img/licence/FreeCulturalWorks_seal_x2.jpg}
		\end{minipage}
	
		
		\begin{center}
			\bfseries\Large
			Selon les conditions suivantes :
		\end{center}
		
		\begin{itemize}
			\item[] \textbf{Attribution} — Vous devez créditer l'Œuvre, intégrer un lien vers la licence et indiquer si des modifications ont été effectuées à l'Oeuvre. Vous devez indiquer ces informations par tous les moyens raisonnables, sans toutefois suggérer que l'Offrant vous soutient ou soutient la façon dont vous avez utilisé son Oeuvre. 
			\item[] \textbf{Pas de restrictions complémentaires} — Vous n'êtes pas autorisé à appliquer des conditions légales ou des mesures techniques qui restreindraient légalement autrui à utiliser l'Oeuvre dans les conditions décrites par la licence.
		\end{itemize}
		
	\end{tcolorbox}
	
%	\begin{center}
%		\bfseries Attribution 4.0 International (CC BY 4.0)
%		\url{https://creativecommons.org/licenses/by/4.0/deed.fr}
%	\end{center}

%	\tiny
%
%	Vous êtes autorisé à : 
%	\begin{itemize}
%		\item \textbf{Partager} — copier, distribuer et communiquer le matériel par tous moyens et sous tous formats
%		\item \textbf{Adapter} — remixer, transformer et créer à partir du matériel
%	\end{itemize}
%	
%	Selon les conditions suivantes : 
%	\begin{itemize}
%		\item \textbf{Attribution} — Vous devez créditer l'Œuvre, intégrer un lien vers la licence et indiquer si des modifications ont été effectuées à l'Oeuvre. Vous devez indiquer ces informations par tous les moyens raisonnables, sans toutefois suggérer que l'Offrant vous soutient ou soutient la façon dont vous avez utilisé son Oeuvre.
%		\item \textbf{Pas d'Utilisation Commerciale} — Vous n'êtes pas autorisé à faire un usage commercial de cette Oeuvre, tout ou partie du matériel la composant. 
%		\item \textbf{Pas de restrictions complémentaires} — Vous n'êtes pas autorisé à appliquer des conditions légales ou des mesures techniques qui restreindraient légalement autrui à utiliser l'Oeuvre dans les conditions décrites par la licence.
%	\end{itemize}

	\end{frame}
}

\settowidth{\leftmargini}{\usebeamertemplate{itemize item}}
\addtolength{\leftmargini}{\labelsep}

\AtBeginDocument{
	\newcolumntype{L}[2]{>{\vbox to #2\bgroup\vfill\flushleft}p{#1}<{\egroup}} 
	
	\begin{frame}[plain]
		\maketitle
	\end{frame}

	\insertlicence
}


% needs etoolbox; to break links after -
\appto\UrlBreaks{\do\-}


\makeatletter
\newcommand{\xRightarrow}[2][]{\ext@arrow 0359\Rightarrowfill@{#1}{#2}}
\makeatother


\usefonttheme{structurebold}
%\usefonttheme{professionalfonts}

\title[TALN : 03- Modèles de langue]%
{Traitement automatique du langage naturel\\Chapitre 03 : Modèles de langue} 

\changegraphpath{../img/modeles/}

\begin{document}
	
\begin{frame}
\frametitle{Traitement automatique du langage naturel}
\framesubtitle{Modèles de langue : Introduction}

\begin{itemize}
	\item Probabilité d'une phrase ($ P(S) = P(w_1, w_2, ..., w_n) $)
	\begin{itemize}
		\item Traduction automatique : \\
		\expword{My tall brother \textrightarrow P(Mon grand frère) \textgreater P(Mon haut frère)}
		\item Correction des fautes grammaticales : \\
		\expword{P(Un objet qu'on peut emporter) \textgreater P(Un objet qu'ont peut emporter)}
		\item Reconnaissance de paroles : \\
		\expword{P(Jeudi matin) \textgreater P(Je dit matin)}
	\end{itemize}
	\item Probabilité d'occurrence d'un mot ($ P(w_n | w_1, \ldots, w_{n-1}) $)
	\begin{itemize}
		\item Auto-complétion : \\
		\expword{P(traitement automatique de l'information) \textgreater P(traitement automatique de l'eau)}
		\item Génération automatique de textes
	\end{itemize}
\end{itemize}

\end{frame}


\begin{frame}
\frametitle{Traitement automatique du langage naturel}
\framesubtitle{Modèles de langue : Plan}

\begin{multicols}{2}
%	\small
\tableofcontents
\end{multicols}
\end{frame}

%===================================================================================
\section{Modèle N-gramme}
%===================================================================================

\begin{frame}
\frametitle{Modèles de langue}
\framesubtitle{Modèle N-gramme}

\begin{block}{Formule des probabilités composées}
	$ P(x_1 \ldots x_n) =  P(x_1) P(x_2 | x_1) P(x_3 | x_1, x_2) \ldots P(x_n | x_1, \ldots, x_{n-1})$
\end{block}


\begin{exampleblock}{Exemple de probabilité d'une phrase}
	$ P(\text{\textit{je travaille à l'ESI}}) =  P(je) P(travaille | je) P(\text{\textit{à}} | \text{\textit{je travaille}}) \ldots P(ESI | \text{\textit{je travaille à l'}})$
	
	$P(ESI | \text{\textit{je travaille à l'}}) = \frac{N(\text{\textit{je travaille à l'ESI}})}{N( \text{\textit{je travaille à l'}})}$ 
	
	Où 
	
	$N$ est le nombre d'occurrence d'une phrase dans un corpus d'entraînement
\end{exampleblock}

\begin{itemize}
	\item Plusieurs phrases possibles
	\item Il faut faut un grand corpus pour estimer cette probabilité (infinité de phrases possibles)
\end{itemize}

\end{frame}

\subsection{Formulation}

\begin{frame}
\frametitle{Modèles de langue : N-gramme}
\framesubtitle{Formulation : Propriété de Markov}

\begin{block}{Propriété de Markov}
	Un état futur ne dépend que de l'état présent. 
	\[%
	P(x_n | x_1,\ldots, x_{n-1}) \approx P(x_n | x_{n-1})
	\]
	Cas général avec $k$ états passés 
	\[%
	P(x_1 \ldots x_n) = P(x_n | x_{n-k}, \ldots, x_{n-1})
	\]
\end{block}

\begin{block}{Estimation de probabilité en utilisant les N-grammes (k grammes)}
	\[
	P(x_n | x_1,\ldots, x_{n-1}) \approx \prod_i P(x_i | x_{i-k}, \ldots, x_{i-1})
	\]
\end{block}

\end{frame}

\begin{frame}
\frametitle{Modèles de langue : N-gramme}
\framesubtitle{Formulation : Quelques modèles}

\begin{itemize}
	\item \optword{Modèle uni-gramme} : $P(w_n | w_1,\ldots, w_{n-1}) \approx P(w_n)$
	\item \optword{Modèle bi-gramme} : $P(w_n | w_1,\ldots, w_{n-1}) \approx P(w_n | w_{n-1})$
	\item \optword{Modèle tri-gramme} :  $P(w_n | w_1,\ldots, w_{n-1}) \approx P(w_n | w_{n-2}, w_{n-1})$
	\item Google Books Ngram Viewer
	\begin{itemize}
		\item \url{https://books.google.com/ngrams}
		\item Modèles pré-traités à partir des livres 
		\item Téléchargement gratuit : \url{https://storage.googleapis.com/books/ngrams/books/datasetsv3.html}
	\end{itemize}
\end{itemize}

\end{frame}

\begin{frame}
\frametitle{Modèles de langue : N-gramme}
\framesubtitle{Formulation : Google Books Ngram Viewer (Humour)}

\hgraphpage{humour-ngram.png}

\end{frame}

\begin{frame}
\frametitle{Modèles de langue : N-gramme}
\framesubtitle{Formulation : Estimation}

\begin{itemize}
	\item en utilisant un corpus d'entraînement avec suffisamment de données
	\item on marque le début et la fin des phrases avec \keyword{\textless s\textgreater} et \keyword{\textless/s\textgreater} (une fois pour les bi-grammes, 2 fois pour les tri-grammes, etc.)
	\item \keyword{l'estimateur du maximum de vraisemblance}
\end{itemize}

\begin{block}{Estimation des probabilités en utilisant le maximum de vraisemblance}
	\[%
	P(w_n | w_{n-k},\ldots, w_{n-1}) = \frac{N(w_{n-k} \ldots w_{n-1} w_n)}{\sum_i N(w_{n-k} \ldots w_{n-1} w_i)}
	= \frac{N(w_{n-k} \ldots w_{n-1} w_n)}{N(w_{n-k} \ldots w_{n-1})}
	\]
	Où $N$ est le nombre d'occurrences des N-grammes dans le corpus
	\[%
	\text{Bi-grammes : } P(w_n | w_{n-1}) = \frac{N(w_{n-1} w_n)}{N(w_{n-1})}
	\]
\end{block}

\end{frame}

\begin{frame}
\frametitle{Modèles de langue : N-gramme}
\framesubtitle{Formulation : Exemple}

\begin{exampleblock}{Exemple d'un corpus d'entraînement}
	\begin{itemize}
		\item \textless s\textgreater un ordianteur peut vous aider \textless/s\textgreater
		\item \textless s\textgreater il veut vous aider \textless/s\textgreater
		\item \textless s\textgreater il veut un ordinateur \textless/s\textgreater
		\item \textless s\textgreater il peut nager \textless/s\textgreater
	\end{itemize}
\end{exampleblock}

\begin{itemize}
	\item $P(peut | il) = \frac{N(il\ peut)}{N(il)} = \frac{1}{3}$
	\item $P(\text{\textit{\textless s\textgreater il peut vous aider \textless/s\textgreater}}) = 
	P(il|\text{\textit{\textless s\textgreater}}) P(peut|il) P(vous|peut) P(aider|vous) P(\text{\textit{\textless/s\textgreater}}|aider) = 
	\frac{3}{4} \frac{1}{3} \frac{1}{2} \frac{2}{2} \frac{2}{2} = \frac{1}{8}
	$
\end{itemize}


\end{frame}

\subsection{Lissage (Smoothing)}

\begin{frame}
\frametitle{Modèles de langue : N-gramme}
\framesubtitle{Lissage (Smoothing)}

\begin{itemize}
	
	\item Si on utilise des petits N-grammes $ \Longrightarrow $ Perte de l'information
	\begin{itemize}
		\item Les langues permettent des dépendances à long terme
		\item \expword{\underline{L'ordinateur} que j'ai utilisé hier à l'ESI pendant la  séance du cours \underline{a planté}}
	\end{itemize}

	\item Si on utilise des grands N-grammes $ \Longrightarrow $ Complexité élevée du modèle
	\begin{itemize}
		\item Il faut un corpus plus grand
		\item Représentation des N-grammes : $|V|^N$ où $V$ est le vocabulaire et $N$ est le nombres de grammes
	\end{itemize}

	\item Problème des N-grammes absents dans le corpus d'entraînement
	\begin{itemize}
		\item $P(\text{\textit{\textless s\textgreater il veut nager \textless/s\textgreater}}) = 
		P(il|\text{\textit{\textless s\textgreater}}) P(veut|il) P(nager|veut)  P(\text{\textit{\textless/s\textgreater}}|nager) = 
		\frac{3}{4} \frac{2}{3} \frac{0}{1} \frac{1}{1} = 0
		$
	\end{itemize}
	\item L'intuition est d'emprunter une petite portion des probabilités des N-grammes existants pour former une probabilité aux N-grammes absents 
\end{itemize}

\end{frame}

\begin{frame}
\frametitle{Modèles de langue : N-gramme}
\framesubtitle{Lissage (Smoothing) : Lidstone}

\begin{block}{Lissage de Lidstone (Bi-grammes comme exemple)}
	\[%
	P(w_n | w_{n-1}) = \frac{N(w_{n-1} w_n) + \alpha}{N(w_{n-1}) + \alpha |V|}
	\]
	Où $V$ est le vocabulaire du modèle
	
	$\alpha = 1$ : \keyword{Lissage de Laplace} 
	
	$\alpha = 0.5$ : \keyword{Loi de Jeffreys-Perks}
\end{block}

\begin{exampleblock}{Exemple : lissage de Laplace}
	\begin{itemize}
		\item Le corpus contient 8 mots différents
		\item Il y a  $8^2 = 64$ bi-grammes possibles
		\item $P(\text{\textit{\textless s\textgreater il veut nager \textless/s\textgreater}}) = 
		P(il|\text{\textit{\textless s\textgreater}}) P(veut|il) P(nager|veut)  P(\text{\textit{\textless/s\textgreater}}|nager) = 
		\frac{3 + 1}{4 + 64} \frac{2 + 1}{3 + 64} \frac{0 + 1}{1 + 64} \frac{1 + 1}{1 + 64} $
	\end{itemize}
\end{exampleblock}

\end{frame}

\begin{frame}
\frametitle{Modèles de langue : N-gramme}
\framesubtitle{Lissage (Smoothing) : Interpolation}

\begin{block}{Interpolation (Tri-grammes comme exemple)}
	\[%
	P_{I}(w_n | w_{n-2} w_{n-1}) = 
	\lambda_3 P(w_n | w_{n-2} w_{n-1}) 
	+ \lambda_2 P(w_n | w_{n-1}) 
	+ \lambda_1 P(w_n) 
	\]
	
	Où $\sum_i \lambda_i = 1$
	
	$\lambda_3$, $\lambda_2$ et $\lambda_1$ sont estimés en utilisant un autre corpus de réglage
\end{block}

%\begin{exampleblock}{Exemple}
%	\begin{itemize}
%		\item Le corpus contient 8 mots différents
%		\item Il y a  $8^2 = 64$ bi-grammes possibles
%		\item $P(\text{\textit{\textless s\textgreater il veut nager \textless/s\textgreater}}) = 
%		P(il|\text{\textit{\textless s\textgreater}}) P(veut|il) P(nager|veut)  P(\text{\textit{\textless/s\textgreater}}|nager) = 
%		\frac{3 + 1}{4 + 64} \frac{2 + 1}{3 + 64} \frac{0 + 1}{1 + 64} \frac{1 + 1}{1 + 64} $
%	\end{itemize}
%\end{exampleblock}

\end{frame}

\begin{frame}
\frametitle{Modèles de langue : N-gramme}
\framesubtitle{Lissage (Smoothing) : Back-off de Katz}

\begin{block}{Back-off de Katz (Tri-grammes comme exemple)}
	\[%
	P_{BO}(w_n | w_{n-2} w_{n-1}) = 
	\begin{cases}
	P^*(w_n | w_{n-2} w_{n-1}) & \text{si } N(w_{n-2} w_{n-1} w_n) > 0 \\
	\alpha(w_{n-2} w_{n-1}) P_{BO}(w_n | w_{n-1}) & \text{sinon}
	\end{cases}
	\]
	
	Où : 
	
	$P^*$ est la probabilité réduite (la réduction sera distribuée sur les probabilités des N-grammes de l'ordre inférieur)
	
	$\alpha$ est une fonction qui distribue la réduction selon le contexte
\end{block}
%
%\begin{exampleblock}{Exemple}
%	\begin{itemize}
%		\item Le corpus contient 8 mots différents
%		\item Il y a  $8^2 = 64$ bi-grammes possibles
%		\item $P(\text{\textit{\textless s\textgreater il veut nager \textless/s\textgreater}}) = 
%		P(il|\text{\textit{\textless s\textgreater}}) P(veut|il) P(nager|veut)  P(\text{\textit{\textless/s\textgreater}}|nager) = 
%		\frac{3 + 1}{4 + 64} \frac{2 + 1}{3 + 64} \frac{0 + 1}{1 + 64} \frac{1 + 1}{1 + 64} $
%	\end{itemize}
%\end{exampleblock}
%
\end{frame}


%===================================================================================
\section{Modèles neuronaux}
%===================================================================================

\begin{frame}
\frametitle{Modèles de langue}
\framesubtitle{Modèles neuronaux}

\begin{itemize}
	\item 
\end{itemize}

\end{frame}

\subsection{Réseau de neurones à propagation avant}

\begin{frame}
\frametitle{Modèles de langue : Modèles neuronaux}
\framesubtitle{Réseau de neurones à propagation avant}

\begin{itemize}
	\item 
\end{itemize}

\end{frame}

\subsection{Réseau de neurones récurrents}

\begin{frame}
\frametitle{Modèles de langue : Modèles neuronaux}
\framesubtitle{Réseau de neurones récurrents}

\begin{itemize}
	\item 
\end{itemize}

\end{frame}

\subsection{Mots hors vocabulaire}

\begin{frame}
\frametitle{Modèles de langue : Modèles neuronaux}
\framesubtitle{Mots hors vocabulaire}

\begin{itemize}
	\item 
\end{itemize}

\end{frame}

%===================================================================================
\section{Évaluation}
%===================================================================================

\begin{frame}
\frametitle{Modèles de langue}
\framesubtitle{Évaluation}

\begin{itemize}
	\item \optword{Évaluation extrinsèque}
	\begin{itemize}
		\item Évaluer le modèle par rapport à une autre tâche : son effet
		\item Exemple, \expword{La qualité de traduction automatique en utilisant ce modèle} 
		\item Évaluation très couteuse
	\end{itemize}
	\item \optword{Évaluation intrinsèque}
	\begin{itemize}
		\item Évaluer le modèle par rapport à sa représentation du langage
		\item Exemple, \expword{Comparer deux modèles en se basant sur leurs capacités de représenter un dataset de test} 
		\item Ne garantit pas une bonne performance du modèle pour une tâche donnée
	\end{itemize}
\end{itemize}

\end{frame}

\subsection{Évaluation extrinsèque}

\begin{frame}
\frametitle{Modèles de langue : Évaluation}
\framesubtitle{Évaluation extrinsèque}

\begin{itemize}
	\item 
\end{itemize}

\end{frame}

\subsection{Perplexité}

\begin{frame}
\frametitle{Modèles de langue : Évaluation}
\framesubtitle{Perplexité}

\begin{itemize}
	\item 
\end{itemize}

\end{frame}

\insertbibliography{TALN03}{*}

\end{document}

